% Written from the artifacts of the 2026-10-20/21 workstation run. Numbers below are read from files in
% reports/experiments/, reports/data_audit/, artifacts/ or reports/verification/ (see paper_claims.csv); a few are
% simple derivations from them (link counts by subtraction, origin counts from split_manifest.json ranges, the 8%
% sampling fraction) and are described as such.
\subsection{What was actually run}
\label{sec:res:run}
All numbers in this section come from a single-day real-data pass of the released pipeline on the authors'
workstation (Windows 11, NVIDIA GeForce RTX 4050 Laptop GPU, SUMO 1.28.0, Ollama with \texttt{qwen2.5:7b}); the
early stages (graph, features, baselines) were re-run during the session and the artifacts cited here are those of
the final chain, with the stage log in \texttt{reports/verification/stage\_log.jsonl} and run ids and hashes in
\texttt{mlflow\_runs.json} and \texttt{evidence\_registry.csv}. Of the 26 Manhattan DOT links, seven on the West
Side Highway corridor (12th Avenue, 11th Avenue and Westside Highway segments) have no valid speed observation
(positive and at most 100\,mph) in any 2026 record and never enter the 5-minute grid
(\texttt{reports/data\_audit/link\_split\_coverage.csv}); of the 19 gridded links, 16 pass the training-coverage rule and form the modelling set, and three of
those (FDR N to the Triborough Bridge, Triborough Bridge westbound, BQE N to the Manhattan Bridge) have no valid
observation inside the test partition and are reported as NOT\_EVALUATED in
\texttt{reports/experiments/link\_test\_support\_operational.csv}. Of 1{,}968{,}303 raw speed records, 978{,}018
carry a zero speed and 59 exceed the plausibility bound; 990{,}226 are valid and the gridded observed fraction is
0.611. The chronological partitions over the 84{,}084 grid steps are 60/15/10/15\,\%: calibration covers 2026-08-08
04:20 to 09-06 08:55\,UTC and the test partition 2026-09-06 09:00 to 2026-10-20 04:00\,UTC (about 44 days spanning 45 calendar dates; 144{,}217,
144{,}213 and 144{,}205 valid targets at 15, 30 and 60 minutes, identical for every model).

To finish inside a fixed compute deadline, every deep model was trained with a recorded budget of 4{,}096 training
windows per epoch, sampled afresh each epoch without replacement from the 50{,}415 training origins (about 8\,\% of
them), with early stopping on the mean validation MAE of a fixed validation subset (patience 5, at most 30 epochs)
and one seed (42). The budget is identical for
the GRU, the temporal Transformer, the proposed model and its ablations, and is stated here because it bounds how far
each model was trained, not because it is a tuned hyper-parameter. Under it the proposed model stopped after 14 epochs
(best epoch 8, 2.3 minutes on the GPU, 158{,}403 parameters), the temporal Transformer after 14 (best 8, 1.8 minutes),
the GRU after 17 (best 11, 6.0 minutes, 41{,}667 parameters) and the no-weather ablation ran the full 30 epochs (best
24).

\subsection{Multi-horizon forecasting}
Table~\ref{tab:forecast} gives the test error on the operational track. Every learned model beats persistence and the
historical average at every horizon, and the three sequence models beat the two gradient-boosted baselines on MAE
at every horizon. The proposed context-gated GATv2--Transformer has the lowest MAE at 30 and 60 minutes (5.33 and
6.15\,mph against 5.34/6.18 for the temporal Transformer and 5.37/6.17 for the GRU), the GRU the lowest at 15
minutes (4.71\,mph against 4.77 for the proposed model), and XGBoost the lowest 60-minute RMSE (8.98 against 9.05).
Relative to persistence the proposed model reduces MAE by 7\,\%, 12\,\% and 16\,\% at the three horizons and raises
$R^2$ from 0.55/0.46/0.30 to 0.69/0.63/0.56. The margins between the three deep models are 0.01 to 0.06\,mph, which
a single seed cannot separate, so this paper does not claim that the proposed architecture forecasts better than a
well-trained GRU or temporal Transformer on this sparse 16-link graph; it claims that it is comparably accurate (within
0.06\,mph MAE of both at every horizon) while carrying the availability gate and the graph layer that the rest of the workflow uses. On the retrospective
track (reanalysis weather plus collisions and permitted events, contextual analysis only) the proposed model reaches
4.67, 5.35 and 6.13\,mph (Table~\ref{tab:forecast} is the operational table; the retrospective rows are in
\texttt{forecast\_metrics.csv}), a small gain at 15 and 60 minutes over the operational 4.77/5.33/6.15, with XGBoost
and LightGBM at 4.89, 5.61--5.65 and 6.30--6.39.

\textit{Ablations.} Table~\ref{tab:ablation} reports the mean MAE over the three horizons and a paired day-block
bootstrap (1{,}000 resamples, the code default, of the 44 test days that contain a valid paired target; 95\,\%
percentile interval) of the difference to the
proposed model (5.419\,mph). Removing the availability gate (plain concatenation) raises MAE by 0.067\,mph
($[0.048, 0.086]$); removing the 311 infrastructure group raises it by 0.014\,mph ($[0.002, 0.026]$); removing the
operational weather forecasts \emph{lowers} it by 0.046\,mph ($[-0.069, -0.021]$). The last is a negative result:
with this data, this budget and the re-forecast weather proxy, the weather features hurt. The no-events and
no-incidents variants are NOT\_APPLICABLE on the operational track, where those groups do not exist, and are listed
as such rather than as ablations. An earlier execution of this stage had resampled 5-minute steps instead of days
(a unit error in the day index) and produced far narrower intervals; it was found during the run, corrected in code,
re-run, and the superseded file is kept with the run's local evidence.

\subsection{Uncertainty calibration and the reliability gate}
Split-conformal half-widths fitted on 102{,}301, 102{,}299 and 102{,}302 calibration residuals at 15, 30 and 60
minutes are 12.1, 13.5 and 14.8\,mph.
On the test partition the 90\,\% intervals reach empirical coverage of 0.909, 0.907 and 0.903 (Table~\ref{tab:unc}),
so the marginal guarantee holds, but the intervals are wide (mean widths 24.3, 26.9 and 29.5\,mph, twice the conformal
quantiles). Conditional coverage is uneven: 0.93 on free-flowing targets (speed ratio
$\ge 0.8$, $n\approx 98{,}500$), 0.93 on moderate ones, and 0.80, 0.79 and 0.76 at the three horizons under
congestion (ratio $< 0.5$, $n\approx 23{,}300$); 0.93--0.94 in rush hours against 0.90 outside them; 0.92--0.93 under
wet weather codes ($n\approx 7{,}900$) against 0.90--0.91 when dry. The reliability gate, whose thresholds are fixed
on calibration and validation data (interval width at most 38.1\,mph, the 80th percentile of validation widths; at
most 34\,\% missing recent inputs; data at most 20 minutes old; at least 200 calibration samples), defers 25.5, 25.7
and 26.1\,\% of test forecasts at the three horizons; the retained forecasts have MAE 3.72, 4.38 and 5.25\,mph against
4.77, 5.33 and 6.15 for all forecasts, and coverage 0.95, 0.94 and 0.93 (\texttt{artifacts/calibration/deferral\_tradeoff.csv}
tabulates the trade-off at six tolerated-missing-input thresholds, 0, 10, 20, 34, 50 and 100\,\%).

\subsection{Explanations of the trained model}
Exact grouped KernelSHAP (all $2^P$ coalitions, $P=5$ or 6 players, 528 or 1{,}040 model calls and 1.7 to 4.2\,s per
explanation on the CPU, additivity gap below $10^{-3}$\,mph) was computed on the trained proposed model for 12 test
points: the six largest predicted 30-minute speed drops, one per link, and six random points, all drawn from the
132{,}355 of 201{,}616 test (origin, link) pairs whose link had a raw observation at every one of the last six input
steps and an observed target (the eligibility rule was introduced after a first pass had selected a point on a link with no raw
observations in the test period). Averaged over the 12 points, the mean absolute attribution is 8.22\,mph for the target
link's own last 30 minutes, 2.04 for its earlier 90 minutes, 1.09 for calendar, 0.52 for graph-neighbour traffic,
0.19 for weather and 0.09 for 311 requests. In the largest alert, the George Washington Bridge lower-level plaza
link at 30 minutes, the model predicts 2.25\,mph against a 26.9\,mph background and an observed 1.24\,mph, with
$-21.3$\,mph attributed to the link's own recent traffic; on FDR N 25th--63rd St it predicts 1.92\,mph (observed 6.83)
with $-13.9$ from recent traffic and $-3.2$ from the earlier window. One selected point is a clear miss: FDR S
Catherine Slip to the Brooklyn Bridge predicted 27.1\,mph where 2.48 was observed, a drop not foreseeable from the
inputs. Shortening the input history confirms the attribution: on a 20{,}665-target test subset the proposed model's
MAE rises from 5.43\,mph with the full two hours to 6.00 (60 min), 6.29 (30), 6.54 (15) and 7.11 (5 min). Removing
all edges touching a link's graph neighbours changes its predictions by 1.34\,mph on average for the four-neighbour
Brooklyn Bridge--FDR N link (MAE 8.57 to 8.58), by 0.91 and 0.87 for FDR S 25th St--Catherine Slip and FDR N
Catherine Slip--25th St (MAE 5.27 to 5.17 and 8.13 to 7.96, i.e.\ slightly better without the edges) and by 0.67 for
the Brooklyn--Battery Tunnel eastbound approach, and by exactly 0.0 for West St S at Spring St, whose only remaining
edges are outgoing (its upstream neighbour is a dead feed), so no message reaches it. The LightGBM TreeSHAP baseline
attributes 65\,\% of its mean absolute attribution at 15 minutes to traffic lags, 16\,\% to calendar and 4\,\% each
to weather and link identity. All of these are model sensitivities, not causal effects.

\subsection{Evidence retrieval}
The approved corpus holds 15 documents (the FHWA 2008 Signal Timing Manual and its chapters 5 and 6, MUTCD 2009
chapters 4D and 4E, nine SUMO documentation pages pinned to v1.28.0 and the project-defined policy), 1{,}633 chunks
embedded with all-MiniLM-L6-v2 into a FAISS inner-product index; 171 chunks exceed the encoder's 256-token window and
are truncated by it (recorded in \texttt{index\_manifest.json}). On the 28 project-labelled queries the index reaches
Recall@8 of 0.818 and MRR of 0.702, and the support-status decision (SUPPORTED against INSUFFICIENT\_EVIDENCE at a
0.35 score threshold) matches the label on 27 of 28 queries (0.964, Table~\ref{tab:ret}). The labels were drafted by
the project and have not been reviewed externally.

\subsection{Advisory workflow}
Table~\ref{tab:adv} lists all 19 policy-by-case runs (REPLAY mode on test-partition forecast origins). The deterministic
gates fired as designed (the three stress cases were run under the full system only, the rainy period under every
policy): the stale-data stress case (inputs aged six steps) ended in DEFER\_DATA,
the wide-interval stress case (forced 99\,mph width) and the observed rainy period, whose intervals of 44--46\,mph exceed
the 38.1\,mph cap, in DEFER\_UNCERTAINTY, and the out-of-scope query in DEFER\_EVIDENCE, each without an LLM call. The
model was called 12 times (one attempt each, 892 to 3{,}604 prompt tokens, 157 to 278 completion tokens, 15 to
126\,s per call with the first, cold call slowest) and returned schema-valid JSON every time
(\texttt{reports/experiments/llm\_calls.csv}). Without retrieval it cited an evidence id that does not exist in all
three cases (for example \texttt{EVIDENCE\_123456}), and the citation check rejected all three; with retrieval, no
invalid citation occurred in nine calls. On the routine rush-hour case (FDR S Catherine Slip--Brooklyn Bridge, a link
without a reviewed DOT-to-SUMO mapping) the constraint layer rejected the candidate under every policy that keeps it,
while the RAG-without-constraints policy passed the same candidate to human review. On the two cases on the mapped
West Street link the model recommended monitoring only, which passes all 13 checks (pedestrian protection
NOT\_EVALUATED) and ends in PENDING\_HUMAN\_REVIEW; because no signal change was proposed, the simulation gate was
never invoked by a real advisory and the SIMULATION\_PASSED path remains exercised only by the software tests.
End-to-end latency is 19 to 30\,s per advisory after the first call.

\subsection{Simulation testbed}
\label{sec:res:sim}
The West Street corridor extracted from OpenStreetMap has 82 corridor edges, 22 signalised intersections and 25
side-street entries; demand is 1{,}200 vehicles per hour and direction on the corridor and 150 per side approach,
both declared assumptions because no automated count matched the corridor, and the three scenarios are a routine
peak, an event surge (demand $\times 1.25$) and an incident that, at the same $\times 1.25$ demand, reduces one
corridor edge to 2.5\,m/s for 20 minutes. Two DOT links map to the corridor with high confidence: a dead West Side feed with no usable speed data
(11th Ave S Gansevoort--West St at Spring St) and the modelled West St S Spring St--Brooklyn--Battery Tunnel portal
link. The scenarios' demand and congestion flags are declared rather than derived from that link's observed or
forecast speeds, so the testbed is a transfer demonstration on a real corridor with declared demand rather than a
digital twin of a forecast link. All 45 matched-seed runs (three
scenarios, fixed-time, actuated and forecast-informed-rule controllers, five seeds, 300\,s warm-up and 3{,}600\,s
evaluation, SUMO 1.28.0, 9{,}645\,s of wall time) completed, and Table~\ref{tab:sim} reports their means. The
result is \emph{inconclusive}, and we report it as such: the declared demand oversaturates the corridor under every
controller. Of 6{,}675 (routine) or 8{,}350 (surge and incident) loaded vehicles per run, only 1{,}003 to 3{,}441
complete their trip within the evaluation window, 3{,}257 to 6{,}979 remain unfinished, and 1{,}091 to 2{,}386
vehicles per run are teleported by SUMO after waiting 300\,s in gridlock. Mean time loss per completed trip is
971.6\,s (fixed-time), 987.5 (actuated) and 971.6 (forecast rule) in the routine peak, where the rule never
activates because congestion is not forecast and therefore reproduces fixed-time exactly; 780.5, 1{,}118.4 and 739.0
in the event surge; and 898.4, 1{,}070.1 and 806.0 in the incident scenario. The forecast rule's lower time loss
in the two stress scenarios is survivorship, not improvement: it completes fewer trips than fixed-time (2{,}014
against 2{,}087 in the surge, 1{,}236 against 1{,}384 in the incident) and teleports a comparable number of vehicles (1{,}834 against 1{,}799 in the surge, 2{,}174 against 2{,}268 in
the incident), while the
actuated controller completes the most trips (2{,}877 and 2{,}240) at the highest time loss per completed trip. Under
gridlock the per-trip metrics measure which vehicles got through rather than how well the corridor was served, so no
controller ranking is claimed, the acceptance rule now counts completed trips and teleports alongside delay, and
calibrating the declared demand against observed counts is the first item of future work. Pedestrian performance is
NOT\_EVALUATED because the scenarios contain no pedestrian demand.

\subsection{Runtime}
Training times above are on the GPU; inference of one forecast origin for all links on the CPU costs 5.7\,ms for the
proposed model, 2.2 for the temporal Transformer, 0.7 for the GRU and about 0.1 for XGBoost; a FAISS query takes
tens of milliseconds after the first (model-loading) call; the full per-link runtime table is
\texttt{reports/experiments/runtime\_metrics.csv}.
